\subsection{有界化}

设 E 是复希尔伯特空间。我们用 s 表示酉映射 \((x,y)\mapsto(-y,x)\)（\(E\oplus E\) 上的）。设 u 是 E 上具有稠密定义域的闭部分算子。部分算子 \(u^{*}u\) 自伴且正（IV，第 241 页，命题 12），故 \(-1\notin\mathrm{Sp}(u^{*}u)\)（IV，第 248 页，命题 17）。我们记 \(\mathrm{W}(u)=(1_{\mathrm{E}}+u^{*}u)^{-1}=-\mathrm{R}(u^{*}u,-1)\)；它是 E 的正的单射自同态。

用 \(p_{1}\) 与 \(p_{2}\) 表示 \(\Gamma_{u}\) 到 E 中的两个典范投影。它们是 \(\mathcal{L}(\Gamma_{u};\mathrm{E})\) 的元素，且有对应关系的等式 \(p_{2}=u\circ p_{1}\)。设 \((j,|p_{1}^{*}|)\) 是自同态 \(p_{1}^{*}\in\mathcal{L}(\mathrm{E};\Gamma_{u})\) 的极分解（I，第 140 页，定义 4），即 \(p_{1}^{*}=|p_{1}^{*}|\circ j\)。映射 j 是从 E 到 \(\Gamma_{u}\) 中的部分等距。

定义 1. --- E 的自同态 \(p_2 \circ j\) 称为 \(u\) 的有界化；我们用 \(b(u)\) 表示这个自同态。

我们有 \(\|b(u)\|\leqslant1\)。

命题 1. --- 我们有下列公式

\[
| p _ {1} ^ {*} | = \mathrm{W} (u) ^ {1 / 2},\tag{1}
\]

\[
b (u) = u \circ \mathrm{W} (u) ^ {1 / 2} = u \circ | p _ {1} ^ {*} |,\tag{2}
\]

\[
1 _ {\mathrm{E}} - b (u) ^ {*} b (u) = \mathrm{W} (u),\tag{3}
\]

\[
b (u) \mathrm{W} (u) = \mathrm{W} (u ^ {*}) b (u), \quad b (u) \mathrm{W} (u) ^ {1 / 2} = \mathrm{W} (u ^ {*}) ^ {1 / 2} b (u),\tag{4}
\]

设 \(x \in E\)。令 \(y = \mathrm{W}(u)(x) \in \mathrm{dom}(u^{*}u)\)。我们有 \(y \in \mathrm{dom}(u)\) 且 \(u(y) \in \mathrm{dom}(u^{*})\)。对形如 \((y_{1}, u(y_{1}))\) 的、\(\Gamma_{u}\) 中的每个元素（其中 \(y_{1} \in \mathrm{dom}(u)\)），计算得

\[
\begin{array}{c} \langle (y, u (y)) | (y _ {1}, u (y _ {1})) \rangle = \langle y | y _ {1} \rangle + \langle u (y) | u (y _ {1}) \rangle \\ = \langle y + u ^ {*} u (y) | y _ {1} \rangle = \langle x | p _ {1} (y _ {1}, u (y _ {1})) \rangle . \end{array}
\]

这就表明 \(p_{1}^{*}(x) = (y, u(y))\)，从而 \(p_{1} \circ p_{1}^{*}(x) = y = \mathrm{W}(u)(x)\)。于是有 \(p_{1} \circ p_{1}^{*} = \mathrm{W}(u)\)，故 \(|p_{1}^{*}| = \mathrm{W}(u)^{1/2}\)，即公式 (1)。特别地，由 I，第 140 页命题 11 的推论有 \(\operatorname{Im}(p_{1}) = \operatorname{Im}(|p_{1}^{*}|)\)，由此得出 \(\operatorname{dom}(u) = \operatorname{Im}(p_{1}) = \operatorname{Im}(\mathrm{W}(u)^{1/2})\)。

于是部分算子 \(u \circ \mathrm{W}(u)^{1/2}\) 以 E 为定义域。关系式 \(p_1 \circ j = (p_1^*)^* \circ j = |p_1^*| = \mathrm{W}(u)^{1/2}\)（I，第 140 页，命题 11，a)）进而蕴含

\[
b (u) = p _ {2} \circ j = u \circ p _ {1} \circ j = u \circ \mathrm{W} (u) ^ {1 / 2},
\]

此即公式 (2)。

由 \(\operatorname{Ker}(j) = \operatorname{Ker}(|p_1^*|)\)（I，第 139 页，命题 10，\(b\)）与公式 (1) 便推出部分等距 \(j: \mathrm{E} \to \Gamma_u\) 是单射，从而是等距的，于是 \(j^* \circ j = 1_{\mathrm{E}}\)。把 j 写成 \(j = (p_1 \circ j, b(u))\)，我们得到

\[
1 _ {\mathrm{E}} = j ^ {*} \circ j = (p _ {1} \circ j) ^ {*} (p _ {1} \circ j) + b (u) ^ {*} b (u) = \mathrm{W} (u) + b (u) ^ {*} b (u),
\]

由此得公式 (3)。

设 \(x \in \mathrm{dom}(u)\)，令 \(y = \mathrm{W}(u)(x)\)。我们有 \(y \in \mathrm{dom}(u)\)，且公式 \(y + u^{*}u(y) = x\) 蕴含 \(u^{*}u(y) \in \mathrm{dom}(u)\)。于是我们有

\[
\begin{array}{r l} & {\mathrm{W} (u ^ {*}) ^ {- 1} (u (y)) = (1 _ {\mathrm{E}} + u u ^ {*}) (u (y))} \\ & {\qquad = u (y) + u u ^ {*} u (y) = u (y + u ^ {*} u (y)) = u (x),} \end{array}
\]

这表明 \(u \circ \mathrm{W}(u)\) 到 u 的定义域上的缩减等于 \(\mathrm{W}(u^{*}) \circ u\)。由于 \(|p_{1}^{*}|\) 的像等于 \(p_{1}\) 的像，从而等于 u 的定义域（I，第 140 页命题 11 的推论），由此得

\[
u \circ \mathrm{W} (u) \circ | p _ {1} ^ {*} | = \mathrm{W} (u ^ {*}) \circ u \circ | p _ {1} ^ {*} | = \mathrm{W} (u ^ {*}) \circ b (u)
\]

（公式 (2)）。此外，\(|p_{1}^{*}| = \mathrm{W}(u)^{1/2}\) 与 \(\mathrm{W}(u)\) 交换，于是我们得到

\[
b (u) \circ \mathrm{W} (u) = u \circ | p _ {1} ^ {*} | \circ \mathrm{W} (u) = u \circ \mathrm{W} (u) \circ | p _ {1} ^ {*} | = \mathrm{W} (u ^ {*}) \circ b (u).
\]

此关系式蕴含 \(b(u) \circ f(\mathrm{W}(u)) = f(\mathrm{W}(u^{*})) \circ b(u)\)，对每个函数 \(f \in \mathcal{C}(\mathbf{R}_{+})\)（I，第 113 页，命题 11），特别地有 \(b(u) \circ \mathrm{W}(u)^{1/2} = \mathrm{W}(u^{*})^{1/2} \circ b(u)\)。这就证明了公式 (4) 并完结了证明。

推论。--- 我们有 \(b(u)^{*} = b(u^{*})\)。

用 \(q_{1}\) 与 \(q_{2}\) 表示投影 \(\Gamma_{u^{*}}\to\mathrm{E}\)，并设 \((k,|q_{1}^{*}|)\) 是 \(q_{1}^{*}\) 的极分解，于是 \(b(u^{*})=q_{2}\circ k\)。对 E 中所有 \(x\) 与 \(y\)，我们有 \(j(x)\in\Gamma_{u}\) 与 \(k(y)\in\Gamma_{u^{*}}\)，且由 IV，第 236 页命题 7，\(\Gamma_{u}\) 与 \(s(\Gamma_{u^{*}})\) 正交，由此得

\[
0 = \langle j (x) | s (k (y)) \rangle = \langle p _ {1} \circ j (x) | - q _ {2} \circ k (y) \rangle + \langle p _ {2} \circ j (x) | q _ {1} \circ k (y) \rangle .
\]

由于 \(p_1 \circ j = |p_1^*|\) 且 \(q_1 \circ k = |q_1^*|\)（I，第 140 页，命题 11，b)），我们得到

\[
\langle b (u) x \mid | q _ {1} ^ {*} | y \rangle = \langle | p _ {1} ^ {*} | x \mid b (u ^ {*}) y \rangle ,
\]

故 \(|q_1^*|b(u) = b(u^*)^* |p_1^*|\)。利用公式 (1) 与 (4)，我们由此得出 \(b(u)|p_1^*| = |q_1^*|b(u) = b(u^*)^* |p_1^*|\)，又因 \(|p_1^*|\) 的像是 \(u\) 的定义域，它在 E 中稠密，故得 \(b(u) = b(u^*)^*\)。

我们用 \(\Omega(\mathrm{E})\) 表示由 \(v \in \mathcal{L}(\mathrm{E})\)（满足 \(1_{\mathrm{E}} - v^{*}v\) 正且单射的）所成的集合。对 \(v \in \Omega(\mathrm{E})\)，埃尔米特自同态 \((1_{\mathrm{E}} - v^{*}v)^{1/2}\) 是单射的，因为它的平方是单射的；我们用 \(\mathrm{B}(v)\) 表示部分算子 \(v \circ ((1_{\mathrm{E}} - v^{*}v)^{1/2})^{-1}\)。

注意 \(1_{\mathrm{E}} - v^{*}v\) 正当且仅当 \(\| v\| \leqslant 1\)，因为 \(\langle x|(1_{\mathrm{E}} - v^{*}v)(x)\rangle = \| x\| ^2 -\| v(x)\| ^2\) 对一切 \(x\in \mathrm{E}\) 成立。

引理 1. --- 子集 \(\Omega(\mathrm{E})\) 在 \(\mathcal{L}(\mathrm{E})\) 中是自伴的。

设 \(v \in \Omega(\mathrm{E})\)。我们有 \(\|v^{*}\| = \|v\| \leqslant 1\)。于是自同态 \(1_{\mathrm{E}} - vv^{*}\) 是正的。它是单射的：若 \(x \in \operatorname{Ker}(1_{\mathrm{E}} - vv^{*})\)，我们有 \(vv^{*}(x) = x\)，故 \(v^{*}(v(v^{*}(x))) = v^{*}(x)\)，于是 \(v^{*}(x) = 0\)（因为 \(1_{\mathrm{E}} - v^{*}v\) 单射），最后 \(x = v(v^{*}(x)) = 0\)。引理由此得证。

命题 2. — 映射 \(u \mapsto b(u)\) 定义了从 E 上闭的稠定部分算子集合到集合 \(\Omega(E)\) 上的一个双射。逆双射由 \(v \mapsto B(v)\) 给出。

关系式 \(1_{\mathrm{E}} - b(u)^{*}b(u) = \mathrm{W}(u)\)（公式 (3)）蕴含 \(b(u)\) 属于 \(\Omega (\mathrm{E})\)，因为 \(\mathrm{W}(u)\) 正且单射。此外，由 \(b(u) = u\circ \mathrm{W}(u)^{1 / 2}\)（公式 (2)），我们得到

\[
u = b (u) \circ (\mathrm{W} (u) ^ {1 / 2}) ^ {- 1} = b (u) \circ ((1 - b (u) ^ {*} b (u)) ^ {1 / 2}) ^ {- 1} = \mathrm{B} (b (u)).
\]

反之，设 \(v \in \Omega(\mathrm{E})\)。令 \(w = (1_{\mathrm{E}} - v^{*}v)^{1/2}\)，并令 \(u = \mathrm{B}(v) = v \circ w^{-1}\)。u 的定义域是 w 的像，由于 w 埃尔米特且单射，它在 E 中稠密（EVT，V，第 41 页，命题 4 (i)）。于是对每个 \(x \in \mathrm{dom}(u)\)，有

\[
\begin{array}{r l} & {\| u (x) \| ^ {2} = \langle (v ^ {*} v) (w ^ {- 1} (x)) | w ^ {- 1} (x) \rangle} \\ & {\qquad = - \langle (1 _ {\mathrm{E}} - v ^ {*} v) (w ^ {- 1} (x)) | w ^ {- 1} (x) \rangle + \| w ^ {- 1} (x) \| ^ {2}} \\ & {\qquad = - \langle w (x) | w ^ {- 1} (x) \rangle + \| w ^ {- 1} (x) \| ^ {2}} \end{array}
\]

因为 \((1_{\mathrm{E}} - v^{*}v) \circ w^{-1}\) 是 w 到 \(w^{-1}\) 的定义域上的缩减。由于 w 自伴，由此得

\[
\| w ^ {- 1} (x) \| ^ {2} = \| x \| ^ {2} + \| v \circ w ^ {- 1} (x) \| ^ {2} = \| x \| ^ {2} + \| u (x) \| ^ {2},
\]

于是由 IV，第 231 页引理 2，部分算子 \(\mathrm{B}(v) = u = v \circ w^{-1}\) 是闭的。

最后我们来证明 \(b(\mathrm{B}(v)) = v\)。由于 \(u = v \circ w^{-1}\)，我们有 \(v = u \circ w\)，因此由公式 (2)，只需验证 \(w = \mathrm{W}(u)^{1/2}\)，甚至只需验证 \(1_{\mathrm{E}} - v^{*}v = \mathrm{W}(u)\)。

自同态 \(v^*\) 属于 \(\Omega(\mathrm{E})\)。由引理 1，令 \(w' = (1_{\mathrm{E}} - vv^{*})^{1/2}\)。它是 E 的单射正自同态，且我们有 \(v \circ w = w' \circ v\)（I，第 113 页，命题 11）。\(u\) 的图像是形如 \((w(x), v(x))\)（\(x \in \mathrm{E}\)）的元素所成的集合。由 IV，第 236 页命题 7，\(u^*\) 的图像是 \(s(\Gamma_u)^{\circ}\)。由于对一切 \(x\) 与 \(y \in \mathrm{E}\)，有

\[
\langle (w ^ {\prime} (x), v ^ {*} (x)) | (- v (y), w (y)) \rangle = - \langle x | w ^ {\prime} v (y) \rangle + \langle x | v w (y) \rangle = 0,
\]

故 \(u^*\) 的图像包含形如 \((w'(x), v^*(x))\)（\(x \in \mathrm{E}\)）的元素，并且此时有 \(u^*(w'(x)) = v^*(x)\)。特别地，\(u^*\) 的定义域包含 \(w'\) 的像。

设 \(x \in E\)，令 \(y = w^{2}(x) = (1_{\mathrm{E}} - v^{*}v)(x)\)，于是 \(x = y + v^{*}v(x)\)。我们有 \(y \in \operatorname{dom}(u)\) 且 \(u(y) = v(w^{-1}(y)) = v(w(x)) = w'(v(x))\)。特别地，\(u(y) \in \operatorname{Im}(w) \subset \operatorname{dom}(u^{*})\)，且 \(u^{*}(u(y)) = v^{*}(v(x))\)。由此得 \(x = y + v^{*}v(x) = y + u^{*}u(y)\)，从而 \(y = \operatorname{W}(u)(x)\)，即

\[
(1 _ {\mathrm{E}} - v ^ {*} v) (x) = \mathrm{W} (u) (x).
\]

我们于是证明了 \(1_{\mathrm{E}} - v^{*}v = \mathrm{W}(u)\)，这正是所要的。

推论。--- 自同态 \(b(u)\) 是埃尔米特的，当且仅当 u 自伴。

这由映射 \(u \mapsto b(u)\) 的单射性（命题 2）以及公式 \(b(u^{*}) = b(u)^{*}\)（命题 1 的推论）得出。

